\section{刚体位姿}
\label{sec:lie-rigid-poses}

一个刚体在空间中的状态由两部分决定：它朝向哪里，以及它的参考点在哪里。前者是上一节的旋转，后者是一个平移。可以证明，保持任意两点距离与空间定向的变换恰好是先旋转、再平移的复合（这里不给出证明）；这样的变换称为\term{刚体变换}（rigid transformation），它描述的状态称为\term{位姿}（pose）。本节建立刚体变换的复合与求逆，说明它在点与方向上的不同作用，最后用它表达坐标系之间的换算。

\paragraph{齐次矩阵}
刚体变换把$\bm x$变成$\bm R\bm x+\bm t$，其中$\bm R\in\mathrm{SO}(3)$、$\bm t\in\mathbb R^3$。它不是线性映射，因为原点被移到$\bm t$。在每个点后面补一个坐标$1$，就能把它写成一次矩阵乘法：
\begin{equation}
 \begin{bmatrix}\bm R\bm x+\bm t\\1\end{bmatrix}=\bm T\begin{bmatrix}\bm x\\1\end{bmatrix},\qquad
 \bm T=\begin{bmatrix}\bm R&\bm t\\\bm0^\top&1\end{bmatrix}.
 \label{eq:lie-homogeneous-pose}
\end{equation}
这样的$4\times4$矩阵全体记为$\mathrm{SE}(3)$，称为\term{三维刚体运动群}（three-dimensional rigid-motion group）。

\paragraph{复合与求逆}
按分块乘法，先作用$\bm T_1$、再作用$\bm T_2$得到
\begin{equation}
 \bm T_2\bm T_1=\begin{bmatrix}\bm R_2\bm R_1&\bm R_2\bm t_1+\bm t_2\\\bm0^\top&1\end{bmatrix},\qquad
 \bm T^{-1}=\begin{bmatrix}\bm R^\top&-\bm R^\top\bm t\\\bm0^\top&1\end{bmatrix}.
 \label{eq:lie-pose-composition}
\end{equation}
乘积仍是同样形状的矩阵，逆矩阵也是，单位元是$\bm I_4$，所以$\mathrm{SE}(3)$是$\mathrm{GL}(4,\mathbb R)$的子群。求逆的公式可以直接解出：由$\bm y=\bm R\bm x+\bm t$得$\bm x=\bm R^\top(\bm y-\bm t)=\bm R^\top\bm y-\bm R^\top\bm t$。

复合的平移部分是$\bm R_2\bm t_1+\bm t_2$，而不是$\bm t_1+\bm t_2$：第一次的平移要随第二次的旋转一起转动。例如两次都绕$\bm e_3$转$\pi/2$，平移分别为$\bm t_1=(1,0,0)^\top$与$\bm t_2=(0,1,0)^\top$。复合的旋转是绕$\bm e_3$转$\pi$，平移是$\bm R_2\bm t_1+\bm t_2=(0,1,0)^\top+(0,1,0)^\top=(0,2,0)^\top$。按另一次序复合，平移变成$\bm R_1\bm t_2+\bm t_1=(-1,0,0)^\top+(1,0,0)^\top=\bm0$，结果不同。

\paragraph{作为Lie群}
$\mathrm{SE}(3)$在$\mathrm{GL}(4,\mathbb R)$中由“左上块在$\mathrm{SO}(3)$中、末行为$(0,0,0,1)$”这些闭条件给出，是闭子群，因而是矩阵Lie群。$(\bm R,\bm t)\mapsto\bm T$是从$\mathrm{SO}(3)\times\mathbb R^3$到$\mathrm{SE}(3)$的微分同胚，所以$\mathrm{SE}(3)$是六维流形：三个旋转自由度加三个平移自由度。但作为群，它不是$\mathrm{SO}(3)$与$\mathbb R^3$的直积：直积的运算会给出$(\bm R_2\bm R_1,\bm t_1+\bm t_2)$，而式~\eqref{eq:lie-pose-composition}中旋转会作用到前一次的平移上。

对经过$\bm I_4$的曲线求导，末行恒为零、左上块的导数在$\mathfrak{so}(3)$中，所以Lie代数是
\[
\mathfrak{se}(3)=\left\{\begin{bmatrix}[\bm\omega]_\times&\bm v\\\bm0^\top&0\end{bmatrix}:\bm\omega,\bm v\in\mathbb R^3\right\}.
\]
用归纳法可得，对$k\geq1$，该矩阵的$k$次幂是$\begin{bmatrix}[\bm\omega]_\times^k&[\bm\omega]_\times^{k-1}\bm v\\\bm0^\top&0\end{bmatrix}$，于是
\[
\exp\begin{bmatrix}[\bm\omega]_\times&\bm v\\\bm0^\top&0\end{bmatrix}
=\begin{bmatrix}\exp([\bm\omega]_\times)&\bm V\bm v\\\bm0^\top&1\end{bmatrix},\qquad
\bm V=\sum_{k\geq0}\frac{[\bm\omega]_\times^k}{(k+1)!}.
\]
$\bm\omega=\bm0$时$\bm V=\bm I$，指数就是平移$\bm v$；$\bm\omega\neq\bm0$时，平移部分是$\bm V\bm v$而不是$\bm v$，单参数子群一般绕某条轴转动、同时沿该轴推进（推进量可以为零），是一种螺旋运动。所以$\mathrm{SE}(3)$上的指数坐标不等于“旋转向量加平移向量”这种分开的记录，使用时须说明采用哪一种参数化。

\paragraph{点与方向}
位姿对点和方向的作用不同。点有位置，平移会移动它；方向只有朝向，例如速度、表面法向或两点之差，平移不应改变它。设$\bm x,\bm y$是两个点，它们的差在变换后为$(\bm R\bm x+\bm t)-(\bm R\bm y+\bm t)=\bm R(\bm x-\bm y)$，平移抵消了。在齐次写法中，把点记为$(\bm x,1)$，把方向记为$(\bm d,0)$，则
\[
\bm T\begin{bmatrix}\bm d\\0\end{bmatrix}=\begin{bmatrix}\bm R\bm d\\0\end{bmatrix},
\]
末坐标$0$恰好让平移列不起作用，两个点之差的末坐标也自然是$1-1=0$。

点与方向的区分在计算中容易出错。把方向向量当成点，用$\bm R\bm d+\bm t$变换，会给方向加上一个与位姿有关的偏移；把点当成方向，则会丢掉平移。若变换是一般的可逆线性映射$\bm A$而不是旋转，法向量要按$\bm A^{-\top}$变换才能保持与切向量的正交；对旋转，$\bm R^{-\top}=\bm R$，法向与其他方向按同一规则变换。

\paragraph{坐标系变换}
同一个点在不同坐标系中有不同的坐标，刚体变换也负责这种换算。设坐标系$B$的原点在坐标系$A$中的坐标为$\bm t_{AB}$，$B$的三个坐标轴在$A$中的方向依次是$\bm R_{AB}$的三列。若一点在$B$中的坐标为$\bm x_B$，它在$A$中的位置是$B$的原点加上沿$B$各轴的位移，即
\[
\bm x_A=\bm R_{AB}\bm x_B+\bm t_{AB},\qquad
\bm T_{AB}=\begin{bmatrix}\bm R_{AB}&\bm t_{AB}\\\bm0^\top&1\end{bmatrix}.
\]
这里$\bm T_{AB}$读作“从$B$的坐标到$A$的坐标”。由此得到两条规则：经过中间坐标系$B$换算，$\bm T_{AC}=\bm T_{AB}\bm T_{BC}$，相邻下标相同才能相乘；反向换算是求逆，$\bm T_{BA}=\bm T_{AB}^{-1}$。

同一个矩阵$\bm T_{AB}$有两种读法：它可以看作把点从一个位置移到另一个位置的主动变换，也可以看作同一个点在两个坐标系之间的坐标换算。两种读法使用同一套矩阵运算，但描述的对象不同；混淆二者，最常见的后果是该用$\bm T$的地方用了$\bm T^{-1}$，或者乘积次序颠倒。

给定一组点在两个坐标系中的对应坐标，要反求$\bm R_{AB}$与$\bm t_{AB}$，这就是由点对应估计刚体变换的问题。它需要拟合目标、正交约束下的求解与可辨识性分析，在第~\ref{chap:spatial-map-estimation}章研究。
